\newpage
\section{Conclusion}
\label{sec:conclusion}


This thesis investigated energy-aware configuration control for locally
deployed RAG systems, asking whether per-query module activation can
reduce energy consumption without degrading answer quality.  We
hypothesised that a circular dependency limits pre-execution routing:
the information needed to select modules is produced by executing those
modules.

We designed a measurement framework that attributes energy, latency and
quality to individual pipeline modules across infrastructure nodes.
Applying it to the CITI KMS system with 1{'}200~HotpotQA queries across
eight configurations (9{'}600 sessions), we established three findings.

First, always-on module activation is measurably inefficient.  The full
pipeline consumes 5.6$\times$ more energy than the minimal pipeline, while
quality improves by less than 0.06~points on either metric.  The cost is
unconditional; the benefit is conditional and query-dependent.

Second, an oracle router saves 75\% energy versus always-C7, but this
is primarily an engineering decision: disabling HyDE, which degrades
quality at 5.3$\times$ the energy cost.  In the non-HyDE space
(\{C0, C2, C3, C6\}), the oracle ceiling drops to 19\% versus
always-C6.  The 75\% quantifies removing a harmful module; the 19\%
quantifies the algorithmic routing limit.

Third, eight routing experiments (five strategies, three explorations)
all converge to a config accuracy of 0.38--0.40, with no statistically
meaningful improvement over the cheapest-configuration baseline.  The
ceiling is signal-limited: pre-execution features cannot capture the
query--document interactions that determine which modules help.  The
information needed to route is produced by the computation routing aims
to skip.

These results yield four contributions:

\begin{enumerate}
    \item \textbf{Contribution 1}: An energy measurement framework for distributed RAG
          systems that attributes cost per module, per node and per
          stage, reusable for any modular RAG pipeline.
    \item \textbf{Contribution 2}: Empirical evidence that always-on module activation
          is inefficient, with energy costs spanning 5.6$\times$ for
          quality differences below 0.06~points on either metric.
    \item \textbf{Contribution 3}: A systematic routing investigation that quantifies
          the signal boundary for pre-execution module prediction.
          Module-activation routing is harder than model routing because
          quality gaps between configurations (2--4~pp) are an order of
          magnitude smaller than between models (10--30~pp).
    \item \textbf{Contribution 4}: Characterisation of the information gap between
          oracle routing (100\% config accuracy) and achievable
          pre-execution routing (39.5\%), defining a concrete research
          target: cheaper access to post-retrieval signals.
\end{enumerate}

The 60.5~percentage-point gap between oracle and achieved performance
is the thesis's most consequential result: a structural property of
module-activation routing with pre-execution features, not a limitation
to overcome within the current paradigm.  Closing this gap requires
post-retrieval information at a cost lower than the computation it
replaces.  Meanwhile, a well-tuned energy penalty~$\lambda$ captures
most achievable benefit without any learned model, reducing the problem
to a policy decision rather than a machine learning task.
